\section{Experimental setup}
\subsection{DeformNeXt-10k construction}
\label{sec:data-construction}
We construct paired shape edits from the textured assets and hierarchical
part annotations of PartNeXt~\cite{wang2025partnext}. A canonical source
$M_0=(V_0,F,\mathcal A_0)$ defines vertex indices, ordered faces, and the
UV/material/corner-attribute authority. Each accepted variant is
$M_{1,k}=(V_{1,k},F,\mathcal A_0)$: positions change while correspondence
and face connectivity/order remain fixed. Shared face indices and barycentric
coordinates therefore identify corresponding source and edited surface
points without nearest-neighbor reassignment.

\paragraph{Donor-guided proposals.}
We select donors from the same category using frozen semantic-role and
structural-compatibility heuristics. Semantic roles can contain multiple
disconnected pieces; matching individual leaf meshes or requiring equal
part counts would unnecessarily restrict shape diversity. Source/donor
identities and configuration hashes are recorded before each run. Eight
original-appearance donor renders personalize an SDXL prior with
GeoStyle~\cite{choi2026geostyle}. The donor is a morphology prior, rather
than target geometry or an input to the retargeting model.

GeoStyle proposes face-Jacobian changes under image-space score distillation.
Our production path adds a category prompt and source assembly constraints
to the reconstruction: selected primary surfaces deform, complementary
physical components follow rigid motions, and source-derived attachment
proxies couple their movement. These contacts are geometric estimates,
not annotated physical joints. The main production configuration uses
800 donor-specific LoRA updates and 300 deformation iterations; donor LoRA
and source preprocessing are reused across pairs. This shortened configuration
was adopted for throughput and is not demonstrated equivalent to the
official 3,600-iteration configuration. Supplementary details give the
other changes and retained early configurations.

\paragraph{Final-pair selection.}
The category-prompt/assembly proposal (PA) and an optional 120-iteration
Jacobian, crease-aware normal, and neighboring-displacement refinement
(balanced JN-L) are separate candidates. Each is assessed anew against
$M_0$. Agreement with the initial proposal or the donor is not an admission
criterion. We check canonical correspondence, finite geometry and
catastrophic failures, meaningful deformation, paired RGB usability,
native-representation survival, and actual training-loader compatibility.
Multi-view model-assisted screening identifies obvious spikes, detached
components and assembly collapse; it is not independent human validation.
Small intersections, gaps and folds are recorded as diagnostics rather
than universally rejected. The dataset is therefore not certified
intersection-free. Both meshes retain one source-defined 0.48 frame and
are rendered with the source materials and matched cameras; edited meshes
are not independently normalized. All approved distinct variants are
retained, with near-identical final geometries merged.

\begin{figure*}[t]
\centering
\includegraphics[width=\textwidth]{figures/deformnext-construction.pdf}
\caption{DeformNeXt construction and approved source/edited pairs. Donor
images guide morphology; source assembly constrains reconstruction, and each
final candidate is screened independently. Examples show archived RGB
(camera 8) and clay (camera 14), with a common foreground-union display crop
and scale for M0/M1 at each camera. Four categories are drawn from a pre-existing
large-edit illustration selection whose membership was checked against the
final snapshot. They illustrate visible changes, not typical quality,
generator yield or a 300-versus-3,600-step comparison.}
\label{fig:deformnext-construction}
\end{figure*}

\paragraph{Scope and splits.}
The final snapshot contains 10,165 distinct deformations from 3,483 source
UIDs across 50 categories, averaging 2.92 variants per source. This is a
deformation count, not 10,000 independent source objects or an estimate of
yield over all PartNeXt assets. After construction, a fixed seed assigns
approximately 90/5/5 of the examples to train/validation/test while keeping
every source UID and source geometry cluster in one split
(Table~\ref{tab:deformnext-splits}). Donor disjointness is not enforced;
the held-out evaluation measures retargeting on admitted edits from this
construction process, not generalization to an unseen deformation generator.

\begin{table}[t]
\centering
\small
\caption{Frozen DeformNeXt splits. Variants of one source and geometry
cluster stay together; category proportions are not forced to be uniform.}
\label{tab:deformnext-splits}
\begin{tabular}{lrrrr}
\toprule
Split & Edits & UIDs & Clusters & Classes \\
\midrule
Train & 9,149 & 3,184 & 1,790 & 50 \\
Validation & 508 & 153 & 132 & 42 \\
Test & 508 & 146 & 134 & 42 \\
\midrule
Total & 10,165 & 3,483 & 2,056 & 50 \\
\bottomrule
\end{tabular}
\end{table}


\subsection{Training and internal oracle-scale evaluation}
\paragraph{Data and initialization.}
We use the frozen splits in Table~\ref{tab:deformnext-splits}.
PBR training excludes ten recorded oversized training examples,
using 9,139; no test asset is excluded by that training limit. Source UID and
source-asset clusters are disjoint across these deformation splits. Overlap
with earlier multi-view pretraining is not established. The earlier mixed385
subset of 80,892 training and 950 validation assets is pretraining lineage,
not the deformation fine-tuning set.

We load the original configurations and complete base+LoRA weights:
SS64/Shape512 step 4,210 and Shape1024/PBR1024 step 16,850. The matching backbone
is VGGT-Omega epoch 7. SS uses 480 updates, global batch 96, LoRA rank/alpha 16,
and denoiser/conditioner learning rates $10^{-5}/10^{-4}$. PBR uses 1,900
updates, global batch 24, and rates $5\times10^{-6}/10^{-4}$. These are
approximately five training passes. Ablations retain paired inputs,
initialization, budget, sampler, and noise. Removed modules/channels are
physically absent, not filled with zeros.

\paragraph{Development and freeze.}
SS diagnostics use DEV128/93 sources and three paired generation noises,
averaged within an example. Surface ablations use the same 128 examples,
one noise, frozen Shape and common Top-1 R PBR. PBR appearance uses DEV32
with identical generated geometry, archived guides, texture noise, actual
rendering triangles, and references. SS selection uses CD/F-Score; PBR
selection uses source LPIPS/DreamSim with complementary CLIP. Cross-SS
appearance is a secondary diagnostic, not a veto of geometric improvement.

Two independent SS training seeds for V1/R/Fixed and two PBR seeds for
vanilla/R are reported separately. The SS confirmation concerns occupancy;
PBR confirmation is a four-example GT-shape target-camera training diagnostic.
Neither is two-seed full-pipeline quality confirmation. Two fresh loads on
VAL36 verify the selected SS checkpoint, archived latent/occupancy arrays,
prefix32, and the final batch of four. Actual weight/config/cache bindings
and the complete evaluation protocol are frozen before TEST inference.

\paragraph{Comparators and generation.}
The earlier internal oracle-scale TEST comparison contains five physical pipelines: original
source-only RGB4; original edited-target RGB1+source RGB3; vanilla V1;
Top-1 R; and selected Fixed SS+Top-1 R PBR. Source-only includes the original
M0 view at target camera8, not just three sources. Edited inputs are M1 RGB
at8 plus M0 RGB at9/10/11. Shape weights are common. Original zero controls
retain original PBR; V1 uses vanilla 1900 PBR and R/selected use Top-1 R 1900.
SS generation uses batch 8, 12 steps, guidance 3 and the prescribed per-asset
noise map. Native downstream sampling configurations are retained.

Table~\ref{tab:final-main} also includes a post-hoc descriptive input
reference: the original Pixal3D single-view SS, Shape512, Shape1024, and PBR
denoiser checkpoints, with the project's calibrated projective DINO
conditioner, receive only the edited M1 target RGB at camera 8. This is a
project-specific 1024-cascade evaluation, not the official Pixal3D inference
stack. It uses the same 508 TEST examples, per-asset seeds, scoring cameras,
native renderer, and original-material M1 references as the five-model
comparison. The pretrained single-view weights retain their 12-step sampler
defaults, including SS guidance 7.5. There is no multi-view or deformation
fine-tuning. It was evaluated after the five-model freeze, so it is not
included in the adjusted comparison family.

\paragraph{Metrics and references.}
Source LPIPS~\cite{zhang2018lpips}, DreamSim~\cite{fu2023dreamsim}, and separate
OpenAI CLIP image cosine~\cite{radford2021clip} average cameras 9/10. Target
PSNR/SSIM~\cite{wang2004ssim}/2D mask IoU use 8. References use original
materials on edited M1 geometry, not undeformed source pixels. Surface CD
is the symmetric mean unsquared Euclidean nearest-point distance, using
20,000 area-weighted samples. F-Score uses threshold 0.02 in the common
source-defined normalized frame~\cite{tatarchenko2019surface}; there is no
per-output ICP alignment. Voxel IoU is diagnostic. Full-frame image scores
are primary; foreground/GT-only crops are supplementary.

FID~\cite{heusel2017fid}/KID~\cite{binkowski2018kid} use the same 1,524 images
per model/reference, cameras 8--10, official Inception 2048 FP32 features,
and torch-fidelity 0.3.0 float64 reductions. KID uses 100 subsets of 1,000 and
seed 42017, retaining signed raw values. Distribution differences are
descriptive; subset variability is not a source-cluster interval. These are
camera-8/9/10 evaluations, not evidence of general novel-view performance.
The RGB1 reference sees only camera 8, so cameras 9/10 are held out for that
row; the four-view pipelines receive all three scoring cameras as inputs.
This internal table uses the project's native mesh renderer and
black-background image compositing. Its scores are not directly comparable
to the unified Blender/Cycles benchmark below; in particular, the
object-specific M1-derived scale is available at generation only here.

GT raster quality control uses one-pixel edge tolerance and a uniform
whole-pixel rounding of its 1\% allowance. The post-TEST correction and
GT-only audit are disclosed in the Supplement; metric formulas are unchanged.

\paragraph{Statistical scope.}
Means weight deformation examples equally. Paired bootstrap resamples source
clusters, retaining within-source dependence; equal-source sensitivity is
also retained. The development downstream family is six contrasts across
eight endpoints, with 48,000 draws and 48-comparison Bonferroni intervals.
Its geometry endpoints guide SS interpretation; this family is not reduced
after observing results. Stage IoU has its separate six-contrast family.
The selected TEST model versus V1/R uses 4,000 draws with seed 20261022,
eight pairwise endpoints and a 16-comparison family across 134 sources.
Intervals condition on frozen checkpoints and generation, not a population
of training seeds. A zero-containing interval does not prove equivalence.

% Protocol frozen before the camera-calibrated, scale-unprovided TEST evaluation.
% Result-dependent claims and tables are guarded separately in results.tex.
\subsection{Unified benchmark protocol}
\paragraph{Inputs and camera calibration.}
The task provides an edited target RGB view and three RGB views of the
undeformed source object, together with their camera poses, horizontal fields
of view, and acquisition distances. We do not provide an object-specific
scale computed from either the source or the edited target mesh to a
generator. The official single-view Pixal3D control receives only the edited
target RGB and its camera calibration; its projective pipeline uses the native
unit-scale default when the scale argument is absent. Our frozen V1 and
selected multi-view methods use the same image-only scale estimator, trained
on the training split from RGB/A and camera inputs. We report the actual
inputs used by each external method; methods without a calibrated-camera
interface retain their native inference path.

\paragraph{Common population and outputs.}
We evaluate the frozen TEST508 split, report each method's success count out
of 508, and compute cross-method ranking only on the intersection of
successful cases. The new official RGB1 control and our two multi-view
methods use a 1024 cascade, whereas the previously frozen external Pixal3D
baseline uses its native 1536 cascade and estimated cameras. Each retains its
frozen stage sampling settings, which are reported. We do not force other methods to have the same
triangle count: all raw meshes are sampled with the common 20,000-point
surface protocol and rendered at 512 pixels using the same Blender/Cycles
scene, materials, and camera records. For the official single-view Pixal3D
control and our two multi-view methods, we bake a 2048-pixel texture on a
render-only mesh with a preferred budget of 250,000 faces. If native mesh
cleanup or fast simplification fails, or the baked mesh exceeds
0.01 world units of 99th-percentile nearest-neighbor discrepancy from the
untouched output, we retry at 500,000 and then one million faces;
exceptionally complex meshes may then use two or five million faces under the
same acceptance rule. We record the
actual budget per example. The untouched 1024
output is always used for geometric metrics; this bake does not alter the
generation cascade. At budgets of one million faces or higher, we pre-simplify with
\texttt{fast-simplification} 0.2.0 and preserve that render mesh's topology
during PBR baking; the same rule is used for the official and multi-view
Pixal3D controls. Every baked mesh must pass the same 0.01 world-unit check.
Other methods retain their native textured exports under
the common renderer's one-million-face render-only cap. The official
single-view Pixal3D and
the edited-target multi-view pipelines are labeled separately because their
image inputs differ.

\paragraph{Metrics and optional source-mesh registration.}
Source preservation uses LPIPS, DreamSim, and CLIP cosine on cameras 9 and 10.
Global naturalness uses FID and KID on the same common-success image pool.
Target alignment uses PSNR, SSIM, and foreground mask IoU at camera 8, plus
symmetric unsquared Chamfer distance and F-score at threshold 0.02 on the
surface samples. The surface frame is normalized by the input source object's
center and extent, with no per-output alignment in the raw results. A second,
paired condition applies a similarity transform fitted by trimmed symmetric ICP
to the available undeformed source mesh \emph{after} generation. Registration
never uses the edited target mesh, depth, mask, or rendering, and the aligned
textured outputs are rerendered before image metrics are recomputed. The
source-aligned condition forms the main comparison; paired raw results show
the effect of registration separately for target-camera PSNR, SSIM, and mask
IoU on the exact same success subset. We mark the extra source-mesh input
explicitly. Confidence intervals resample source clusters, retaining all
edited examples from a sampled source together. Neither TEST labels nor TEST
metrics choose a scale estimator, model checkpoint, or ICP hyperparameter.

% Development evidence only; do not treat the small native-render PBR panel as
% the final common Cycles evaluation or promote old M1-scale tables here.
% Evidence: docs/superpowers/reports/2026-10-28-scale-ungiven-VAL31-PBR8-paired.json
\paragraph{Scale-unprovided development gate.}
We first hold the image-only scale estimator and the downstream Shape512,
Shape1024, and vanilla PBR weights fixed while changing only the SS stage.
On 31 source-disjoint validation edits, the V1, Retarget, and Selected
variants obtain CD $0.012610/0.012435/0.012367$ and F@0.02
$0.848527/0.850373/0.853088$, respectively. The paired F-score changes
relative to V1 are $+0.001846$ and $+0.004561$; source-cluster bootstrap
95\% intervals are $[-0.007848,+0.011855]$ and
$[-0.005920,+0.016321]$. The intervals include zero, so this panel does not
establish a geometric gain from either module. Its SS occupancy macro-IoU
is $0.443674/0.444928/0.442446$; stage occupancy and decoded mesh quality
are reported separately.

For appearance isolation, we use the same eight image-scale V1 meshes,
initial PBR noise, and native render references for vanilla and Top-1 R PBR.
Their source-camera LPIPS is $0.070836/0.070935$, and target-camera PSNR is
$24.5466/24.5072$ dB. The paired Top-1-minus-vanilla LPIPS and PSNR
intervals are $[-0.000208,+0.000480]$ and $[-0.1018,+0.0174]$ dB.
This small native-render panel does not establish appearance improvement;
the final Blender/Cycles comparison and source-M0 postprocessing remain
separate evaluations.

